\begin{lemma}[A constant bound for the reflected kernel integrand near $r = 1$]
Fix $n \geq 1$ and $j \in \{1,\ldots,n\}$.  For $x,y \in \mathbb{R}^n$ and $r \in (0,1)$ put
\[
  \Psi_{x,y}(r)
    = \frac{y_j - r x_j}{(1-r^2)^{n/2+1}\sqrt{-\log r}}
      \exp\Bigl(-\frac{|x-ry|^2}{1-r^2}\Bigr) ,
\]
the \emph{reflected} integrand of the raw Gaussian Riesz kernel: it is obtained from the
kernel's own $r$-integrand by interchanging the roles of $x$ and $y$ inside the Gaussian
exponent, and it is the quantity the splitting identity
$|rx-y|^2/(1-r^2) = (|y|^2-|x|^2) + |x-ry|^2/(1-r^2)$ produces.  Then for all $R > 0$ and
$\rho > 0$ there is a constant $C > 0$, depending only on $n$, $j$, $R$ and $\rho$, such
that
\[
  \bigl|\Psi_{x,y}(r)\bigr| \leq C
  \qquad\text{whenever } |x| \leq R,\quad |x-y| \geq \rho,\quad \tfrac12 \leq r < 1 .
\]

The range $\tfrac12 \leq r < 1$ is the one that matters: as $r \uparrow 1$ both
$(1-r^2)^{-n/2-1}$ and $(-\log r)^{-1/2}$ blow up, and the assertion is that the Gaussian
factor more than compensates, uniformly in $x$ and $y$ in the stated range.  That is where
the hypothesis $|x-y| \geq \rho$ is used, and it is the only place it is needed: at $r = 1$
the exponent is $-|x-y|^2/0^+$, so the Gaussian decay is available exactly because $x$ and
$y$ are separated.  The hypothesis $|x| \leq R$ enters through the numerator and through
the competition between $|x|$ and $|y|$.

The bound is a \emph{pointwise} constant bound, not a bound on
$\int_{1/2}^1|\Psi_{x,y}|$, and it is deliberately so: the route below produces the
pointwise statement with no extra work, and the consumer needs a pointwise majorant in
order to get integrability on $(0,1)$ at the same time as the size bound.  Integrating the
constant over a set of measure $\tfrac12$ recovers the integral form at once.  No
uniformity in $n$, $R$ or $\rho$ is claimed, and nothing here feeds an endpoint constant.
\end{lemma}
\begin{proof}
Fix $n \geq 1$ and $j$.  Let $R,\rho > 0$ be given and put
\[
  R' = \max(R,1) , \qquad \rho' = \min(\rho,1) , \qquad
  \tau_0 = \min\Bigl(\tfrac12,\ \frac{\rho'}{8R'}\Bigr) .
\]
Then $1 \leq R'$ and $R \leq R'$; and $0 < \rho' \leq 1$ with $\rho' \leq \rho$.  In
particular $0 < R'$, so $\rho'/(8R') > 0$, whence $\tau_0 > 0$ as the minimum of two
positive numbers, while $\tau_0 \leq \tfrac12$ and $\tau_0 \leq \rho'/(8R')$ hold by the
definition of the minimum.  So $R'$, $\rho'$ and $\tau_0$ satisfy the hypotheses of
\noderef{GaussRieszKernelSingularRangeSliceBound}, which therefore supplies a constant
$C_2 > 0$ such that
\[
  \bigl|\Psi_{x,y}(r)\bigr| \leq C_2
  \qquad\text{whenever } |x| \leq R',\quad \rho' \leq |x-y|,\quad \tfrac12 \leq r < 1 .
\]
It remains to see that this constant works for $(R,\rho)$.  Any pair $(x,y)$ with
$|x| \leq R$ satisfies $|x| \leq R'$ because $R \leq R'$, and any pair with
$|x-y| \geq \rho$ satisfies $|x-y| \geq \rho'$ because $\rho' \leq \rho$.  Hence
$C = C_2$ has the required property, and it depends only on $n$, $j$, $R'$ and $\rho'$,
hence only on $n$, $j$, $R$ and $\rho$.
\end{proof}
